% !TEX TS-program = pdflatex
% !TEX encoding = UTF-8 Unicode
% How to contribute: a page of the wiki "Enhanced dissipation for passive scalars", generated from the wiki on 10 November 2026.
% It compiles as it is with pdflatex (Overleaf, TeXShop, TeXstudio, ...); typeset twice for the citations.
% To propose a change: keep the preamble and the line \title{...}; keep only the blocks you change
% (a \section{...} with its content, or a statement or proof with its \label{...}), change them, and post
% the file in a new topic of the Zulip channel of the part. The rules: page "How to contribute" of the wiki.
\documentclass[11pt]{article}
\usepackage[utf8]{inputenc}
\usepackage[T1]{fontenc}
\usepackage{lmodern}
\usepackage[a4paper,margin=2.5cm]{geometry}
\usepackage{amsmath,amssymb,amsthm}
\usepackage{graphicx,tabularx,listings}
\usepackage[skins,breakable]{tcolorbox}
\usepackage[round]{natbib}
\usepackage[hidelinks]{hyperref}
\setlength{\emergencystretch}{2em}

\newtheorem{theorem}{Theorem}
\newtheorem{proposition}[theorem]{Proposition}
\newtheorem{lemma}[theorem]{Lemma}
\newtheorem{corollary}[theorem]{Corollary}
\newtheorem{conjecture}[theorem]{Conjecture}
\newtheorem*{claim}{Claim}
\theoremstyle{definition}
\newtheorem{definition}[theorem]{Definition}
\newtheorem{example}[theorem]{Example}
\theoremstyle{remark}
\newtheorem{remark}[theorem]{Remark}

%% The commands of the wiki. This preamble is the same for every page: do not change it, and do not
%% define commands of your own (the wiki has none).
\makeatletter
% after a statement: \begin{idea}, \begin{architecture}, \begin{note}[Title] (and \begin{proof}[Title])
\newcommand{\ed@runin}[1]{\par\topsep6\p@\@plus6\p@\relax\trivlist\item[\hskip\labelsep\itshape#1\@addpunct{.}]\ignorespaces}
\newenvironment{idea}[1][Idea]{\ed@runin{#1}}{\endtrivlist\@endpefalse}
\newenvironment{architecture}[1][Proof architecture]{\ed@runin{#1}}{\endtrivlist\@endpefalse}
% \begin{note}[Title]: an aside, in a dashed box; its title in italics at the start of the text
\tcbset{ed@note/.style={enhanced jigsaw,breakable,frame hidden,interior hidden,arc=2pt,borderline={0.5pt}{0pt}{black!50,dash pattern=on 3pt off 2pt},
  left=3mm,right=3mm,top=2mm,bottom=2mm,before skip=3mm plus 1mm,after skip=3mm plus 1mm}}
\newenvironment{note}[1][Note]{\begin{tcolorbox}[ed@note]\noindent{\itshape#1\@addpunct{.}}\enspace\ignorespaces}{\end{tcolorbox}}
% text still to be written; labels of the text; the Lean banner of the statement above
\newcommand{\placeholder}[1]{\if@inlabel\else\par\noindent\fi\textit{[#1]}\@ifnextchar\end{}{\par}}
\newcommand{\ed@label}[1]{\mbox{\footnotesize[#1]}}
\newcommand{\demo}[1]{\ed@label{demo text}}
\newcommand{\heuristic}[1]{\ed@label{heuristic}}
\newcommand{\numerics}[1]{\ed@label{numerics}}
\newcommand{\published}[1]{\ed@label{published}}
\newcommand{\preprint}[1]{\ed@label{preprint, to be verified}}
\newcommand{\raw}[1]{\ed@label{raw, unchecked}}
\def\ed@formalized{formalized}\def\ed@image{image}\def\ed@video{video}\def\ed@diagram{diagram}
\newcommand{\lean}[2][]{\par{\noindent\footnotesize\def\ed@a{#2}Lean: \ifx\ed@a\ed@formalized formalized\ifx\relax#1\relax\else\
  (\href{#1}{certificate})\fi\else not yet formalized\fi.\par}}
% the contributors of a section, most recent first (chosen by the editors), in small type under its title
\newcommand{\contributors}[1]{\par{\noindent\footnotesize\itshape Contributors: #1\par}\smallskip}
% the two lines under the title: the page's team and status, its Zulip topic (default: the title)
\def\ed@title{}\let\ed@maketitle\maketitle\renewcommand{\maketitle}{\let\ed@title\@title\ed@maketitle}
\newcommand{\pagestatus}[1]{{\par\centering\small#1\par}}
\newcommand{\zulip}[2][]{{\par\centering\small Discuss on Zulip: \##2, topic
  ``\ifx\relax#1\relax\ed@title\else#1\fi''.\par}\medskip}
% links to other pages of the wiki, anchors, wiki code shown as text
\newcommand{\wikilink}[2]{\ifx\relax#2\relax#1\else#2\fi}
\newcommand{\wikianchor}[1]{\phantomsection\label{#1}}
\newcommand{\wikipage}[1]{}% the first line of a page file of the whole-wiki project (pages/<page>.tex): the page's title
\newcommand{\nowikicode}[1]{\texttt{\frenchspacing#1}}
\newcommand{\nowiki}[1]{#1}
\newcommand{\wikitemplate}[1]{\texttt{\{\{#1\}\}}}
\newcommand{\wikidiv}[2]{\par#2\par}
% \begin{quote}: a sentence set apart, such as the central question of the Overview; centred, in italics
\renewenvironment{quote}{\list{}{\rightmargin\leftmargin}\item\relax\centering\itshape}{\endlist}
% pictures, videos, diagrams: a grey panel with the picture and its caption. The picture (for a video, its poster image)
% is shown when its file, or a .pdf or .png file of the same name, is next to this file, in figures/ or in ../media
\tcbset{ed@media/.style={enhanced,colback=black!4,frame hidden,arc=2pt,left=3mm,right=3mm,top=2.5mm,bottom=2mm,
  before skip=3mm plus 1mm,after skip=3mm plus 1mm}}
\def\ed@poster{}\def\ed@record{}\def\ed@svg{svg}
\define@key{ed@media}{poster}{\def\ed@poster{#1}}\define@key{ed@media}{record}{\def\ed@record{#1}}\define@key{ed@media}{id}{}
\newcommand{\ed@find}[1]{\IfFileExists{#1}{\edef\ed@found{#1}}{\IfFileExists{figures/#1}{\edef\ed@found{figures/#1}}{%
  \IfFileExists{../media/#1}{\edef\ed@found{../media/#1}}{}}}}
\newcommand{\ed@picturefile}[1]{\def\ed@found{}\edef\ed@b{#1}\ifx\ed@b\@empty\else\filename@parse{\ed@b}\edef\ed@base{\filename@base}%
  \edef\ed@ext{\filename@ext}\ifx\ed@ext\ed@svg\else\ed@find{\ed@b}\fi
  \ifx\ed@found\@empty\ed@find{\ed@base.pdf}\fi\ifx\ed@found\@empty\ed@find{\ed@base.png}\fi\fi}
\newcommand{\ed@panel}[2]{\begin{tcolorbox}[ed@media,#1]\ifx\ed@found\@empty\else{\centering\includegraphics[width=\linewidth,
  height=0.35\textheight,keepaspectratio]{\ed@found}\par}\smallskip\fi{\small\textbf{\ifx\ed@a\ed@video Video\else\ifx\ed@a\ed@diagram
  Diagram\else Figure\fi\fi.}\enspace#2\par}{\footnotesize\itshape\ifx\ed@a\ed@video The video plays in the online edition of this page.\fi
  \ifx\ed@record\@empty\else\ifx\ed@a\ed@video\ \fi How it was made: page Numerical methods, record \texttt{\ed@record}.\fi\par}
  \end{tcolorbox}}
% in the text, the panel is a figure (it floats to where it fits); inside a box (a note, a table), it stays in place
\renewcommand{\topfraction}{0.85}\renewcommand{\bottomfraction}{0.6}\renewcommand{\textfraction}{0.1}\renewcommand{\floatpagefraction}{0.8}
\newcommand{\media}[4][]{\par\def\ed@poster{}\def\ed@record{}\setkeys{ed@media}{#1}\def\ed@a{#2}%
  \ifx\ed@a\ed@video\ed@picturefile{\ed@poster}\else\ed@picturefile{#3}\fi
  \ifinner\ed@panel{}{#4}\else\begin{figure}[htbp]\ed@panel{before skip=0pt,after skip=0pt}{#4}\end{figure}\fi}
% tables (\lv{label}{text}; \header{a & b}, \row{a & b}) and code
\newenvironment{levels}[1][]{\par\smallskip\noindent\tabularx{\linewidth}{@{}>{\bfseries\raggedright\arraybackslash}p{0.2\linewidth}X@{}}}{\endtabularx\par\smallskip}
\newcommand{\lv}[2]{#1 & #2\\[3pt]}
\newenvironment{wikitable}[1]{\par\smallskip\noindent\tabularx{\linewidth}{|*{#1}{>{\raggedright\arraybackslash}X|}}\hline}{\endtabularx\par\smallskip}
\newcommand{\header}[1]{#1\\\hline\hline}
\newcommand{\row}[1]{#1\\\hline}
\lstnewenvironment{wikicode}[1][]{\lstset{basicstyle=\ttfamily\footnotesize,breaklines=true,columns=fullflexible,
  keepspaces=true,showstringspaces=false,frame=single}\ifx\relax#1\relax\else\lstset{language=#1}\fi}{}
% the heading of the reference list is the section References of the page
\renewcommand{\bibsection}{}
% characters of the wiki's text that LaTeX does not know
\DeclareUnicodeCharacter{2192}{\ensuremath{\rightarrow}}\DeclareUnicodeCharacter{203A}{\ensuremath{\rangle}}
\DeclareUnicodeCharacter{2713}{\ensuremath{\checkmark}}\DeclareUnicodeCharacter{25CB}{\ensuremath{\circ}}
\DeclareUnicodeCharacter{2248}{\ensuremath{\approx}}\DeclareUnicodeCharacter{2264}{\ensuremath{\leq}}
\DeclareUnicodeCharacter{2265}{\ensuremath{\geq}}\DeclareUnicodeCharacter{2212}{\ensuremath{-}}
\makeatother

\begin{document}
\title{How to contribute}
\author{Wiki \emph{Enhanced dissipation for passive scalars}}
\date{10 November 2026}
\maketitle
\pagestatus{Page team: Michele Dolce · Status: proposal, to be agreed with the editorial board · Text: \demo{}}
\zulip{contributing}

This wiki is written from the bottom up, the way a journal works: members of the community propose changes on the talk page of the page they change, proposals are discussed and reviewed there in the open, an associate editor decides, and accepted changes are written into the wiki. The roles, and who holds them, are on \wikilink{Contributors}{}. \demo{}

\section{Roles}
\begin{itemize}
  \item \textbf{Editorial board}: the scope and the structure of the wiki; appoints the associate editors; settles discussions that end without consensus.
  \item \textbf{Associate editors}: one or more per part, following the talk pages of its pages and its Zulip channel. They run the discussion of each proposal, ask members of the community for a review when it helps, and decide.
  \item \textbf{Community}: propose, review and discuss. Every member can endorse a new member. Editors can remove a member whose participation disturbs the work rather than helping it.
\end{itemize}
The first line of each page names its team. As the wiki grows, busy parts and pages get more associate editors; the roles stay the same.
On the wiki, everyone with a role is a \emph{member} (reads every page and writes on talk pages); the editorial board and the associate editors are also \emph{editors} (edit articles and templates, and write accepted proposals into the pages); administrators give these rights.

\section{Proposing a change}
Every change to a page, from a typo to a new proof, is a \emph{proposal}: a topic on the talk page of the page (tab \emph{Discussion}), which members of the wiki open and discuss. There are two ways to open one.
\begin{itemize}
  \item \emph{One paragraph, a figure, the central question, or a part still to be written} (a typo, a sentence, a formula, a shorter argument, a caption, a shorter video, a missing proof): in the \href{https://endipa.github.io}{edition}, the pencil button next to it shows its TeX or, for a part still to be written, an empty box (the link \emph{Write this part} opens it too). Change the TeX (the buttons of the box add a remark, a lemma with its proof or a displayed formula after the text), write the missing text, or describe the change you ask for (for a picture or a video, for instance), say why, and open the proposal: the wiki shows a new topic on the talk page, already filled in, to be saved. A text too long for a link is copied by the edition, to be pasted into the topic in the wiki.
  \item \emph{Larger changes} (a section, a statement, a proof, a new section): the link \emph{new proposal} under the title of the page in the wiki (\emph{Propose a change} in the edition) opens a new topic with a short form. Download the TeX source of the page: link \emph{TeX source} under the title of every page of the edition (with the figures, in a zip, when the page has some). It is a standard LaTeX article that compiles as it is. The TeX of the whole wiki (\emph{TeX of the whole wiki}, on the Overview) is one LaTeX project, for working offline, in which every page is a file of the folder \emph{pages}; such a file can be proposed in the same way. Change what you want and check that it compiles; then paste in the form the blocks you changed (a section, or a statement or a proof with its label), or the whole file: only what changes is applied, and there is no limit of length. Say what the proposal changes, why, and which sources it uses; for a long proposal, attach its PDF (\wikilink{Special:Upload}{upload it} and link it in the form).
\end{itemize}
Those who are not yet members of the wiki write in the Zulip channel \emph{contributing}, naming the page.

\subsection{How a proposal is decided}
\begin{enumerate}
  \item The community discusses and reviews the proposal in its topic, with the buttons \emph{Support}, \emph{Oppose} and \emph{Comment} of its box, which add a signed reply (the box counts the supports and the oppositions); replies can also be written by hand, signed with \nowikicode{\textasciitilde{}\textasciitilde{}\textasciitilde{}\textasciitilde{}}. The author may revise the proposed text in place; the history of the talk page keeps every version.
  \item When the discussion has settled, an associate editor of the part decides, with the buttons \emph{Accept} or \emph{Decline} in the box at the top of the topic (shown to editors): the box changes status, and the reason they give is added as a signed reply. There is no deadline. If the discussion ends without consensus, the editorial board decides.
  \item An editor writes the accepted proposal into the page with a script. The script converts the TeX to the wiki's format, refuses a paragraph that has changed since the proposal was made, shows exactly what changes on the page, and flags new citations and links. The edit summary links the topic and credits its author, and the box of the topic becomes \emph{applied}, with a link to the change. Decided topics, applied or declined, are moved to the archive of the talk page (its subpage \emph{/Archive}), so that the talk page shows only the proposals under discussion or waiting to be applied.
\end{enumerate}
Small corrections (a typo, a broken link, a missing reference) go the same way, and are usually accepted without discussion. Open and accepted proposals are listed in \wikilink{:Category:Open proposals}{} and \wikilink{:Category:Accepted proposals}{}.
A proposal changes only the blocks it contains. A paper that is not yet on \wikilink{Literature}{} comes with the proposal that cites it (see \emph{Citations} in \wikilink{\#The TeX format}{The TeX format}), and is added to the reference catalogue when the proposal is applied.

\section{Credits}
Under the title of a section, a line in small type names the people who contributed to it, the most recent first: \nowikicode{\{\{Contributed by|\allowbreak{}Name Surname, Name Surname\}\allowbreak{}\}} on the line after the heading (in TeX, \nowikicode{\textbackslash{}contributors\{Name Surname, Name Surname\}} after the section title). A relevant contribution cannot be measured objectively: the editors decide, at their discretion, who is listed, and add a new name at the beginning of the list. A long list shows the three most recent names, and a link opens the others. The history of each page keeps the record of every change and of its author.

\section{The TeX format}
The TeX source of a page is a standard LaTeX article: one file that compiles with pdflatex and nothing else (on Overleaf, or with TeXShop, TeXstudio, ...; typeset twice for the citations). Its preamble is the same for every page and is never changed; a contribution defines no macros, since the wiki has none. Sections are numbered in the PDF only; the website numbers the results within each page (Theorem 1, Lemma 2, ...).
\begin{itemize}
  \item Statements: theorem, proposition, lemma, corollary, conjecture, definition, example, remark, claim. The optional argument is the name, the label is the anchor: \nowikicode{\textbackslash{}begin\{theorem\}\allowbreak{}[name]\textbackslash{}label\{thm-name\}}. Labels start with thm-, prp-, lem-, cor-, def-, clm-, rem- (remarks), cnj- or exm-; a statement sent without a label gets one made from its name or its first words.
  \item After a statement: the environments idea and architecture, and proofs with an optional title and a label, as in \nowikicode{\textbackslash{}begin\{proof\}\allowbreak{}[Proof by hypocoercivity]\textbackslash{}label\{proof-name\}}. A claim and its proof sit inside the proof that uses them.
  \item Formulas: \nowikicode{\$...\$}, \nowikicode{\textbackslash{}[ ... \textbackslash{}]} and align*; no numbered equations and no labels on equations; none of the commands that the wiki's renderer rejects (the script lists them).
  \item Citations: \nowikicode{\textbackslash{}citet\{Key\}} gives \emph{Author (year)} and \nowikicode{\textbackslash{}citep\{Key\}} gives \emph{(Author, year)}, with the keys of \wikilink{Literature}{}; several papers at once: \nowikicode{\textbackslash{}citep\{Key1,Key2,Key3\}}. The reference list at the end of each page is generated from the reference catalogue and contains only the papers that the page cites: a new citation is added to it when the proposal is applied. A paper that is not yet on \wikilink{Literature}{}: give its full reference in the proposal, as \nowikicode{\textbackslash{}bibitem\{NewKey\}\allowbreak{} Authors, title, journal (year)} in a \texttt{thebibliography} at the end of the TeX, with its DOI or arXiv number (from the pencil button: in the reason field); it is added to the catalogue and then to the page.
  \item Links: on the same page, \nowikicode{\textbackslash{}hyperref[thm-name]\{the theorem above\}} rather than \nowikicode{\textbackslash{}ref}, since the wiki has no numbers; to another page, \nowikicode{\textbackslash{}wikilink\{Page\textbackslash{}\#\allowbreak{}anchor\}\allowbreak{}\{text\}}.
  \item Labels: \nowikicode{\textbackslash{}heuristic\{\}}, \nowikicode{\textbackslash{}numerics\{\}}, \nowikicode{\textbackslash{}published\{\}}, \nowikicode{\textbackslash{}preprint\{\}}, \nowikicode{\textbackslash{}raw\{\}}.
  \item A new block goes after an existing one: put the line \nowikicode{\%\% after: Setting} (a section title or a label) just before it.
\end{itemize}

\section{Anatomy of a result}
A result is a statement followed by its levels of detail, each folded until the reader opens it:
\begin{levels}
\lv{Statement}{in the simplest form, constants hidden}
\lv{Idea}{the key mechanism in a few lines}
\lv{Architecture}{the proof as a short chain of steps, later with a diagram}
\lv{Proof}{the complete proof: first its sources (the papers it follows, cited by their key), then the steps; a statement used only in this proof is a \emph{claim}, with its own folded proof inside}
\end{levels}
A statement that several proofs use, or that holds in more generality than one result, goes to the \emph{Toolkit} (\wikilink{Standard analysis}{} or \wikilink{Enhanced-dissipation tools}{}) and is linked from the proofs. A short proof is written in one piece. The specimen below shows the boxes at work; open the folds to see the nesting.

\subsection{Specimen}
\begin{theorem}[specimen]\label{thm-specimen}
The statement, in its simplest form.
\end{theorem}
\begin{idea}
The key mechanism in a few lines, and the heuristic scaling, labelled \heuristic{}.
\end{idea}
\begin{architecture}
\begin{enumerate}
  \item Step 1 reduces the theorem to a claim.
  \item Step 2 proves the claim.
  \item Step 3 concludes with a lemma of the toolkit.
\end{enumerate}
\end{architecture}
\begin{proof}\label{proof-specimen}
\emph{Sources.} This proof follows ...

\textbf{Step 1.} The reduction.
\begin{claim}[specimen]\label{clm-specimen}
A statement used only in this proof.
\end{claim}
\begin{proof}[Proof of the claim]\label{proof-clm-specimen}
The proof of the claim, folded inside the proof that uses it.
\end{proof}
\textbf{Step 2.} The conclusion, with a lemma of the \wikilink{Standard analysis}{toolkit}.
\end{proof}

\section{Labels}
\begin{itemize}
  \item \demo{}: text drafted by Claude (AI), until a member of the page team replaces or approves it.
  \item \heuristic{}: a non-rigorous explanation, kept apart from what is proved.
  \item \numerics{}: a computation, with its record on \wikilink{Numerical methods}{}.
  \item \published{}, \preprint{}, \raw{}: the status of an entry on \wikilink{Related results}{} and \wikilink{Raw data}{}: refereed; posted but not refereed; neither posted nor checked by a human (an AI-generated proof, a numerical observation, a sketch).
  \item Grey quote boxes: AI-generated text, quoted verbatim with its source and date.
\end{itemize}

\section{Pictures, videos and numerics}
A picture, a video or a diagram is one line in the wiki, the template \wikilink{Template:Media}{Media}; in TeX, \nowikicode{\textbackslash{}media\{video\}\allowbreak{}\{file\}\allowbreak{}\{caption\}}. The files are kept in the folder \texttt{wiki/media} of the project; the edition shows them, the wiki shows a box with the caption. Every computed picture or video has a record on \wikilink{Numerical methods}{}, with the method, the parameters, the checks and the code.

\section{Discussion and conduct}
Proposals are discussed on the talk page of the page they change, so that the discussion stays next to the text and in its history. Zulip is the forum for the mathematics: questions and clarifications, ideas for new results and how they fit the established theory, and larger changes to the structure of the wiki, discussed there before they become proposals. It has one channel per part and one topic per question or idea; the link under the title of each page opens the topic named after the page, for questions about it.
\begin{wikitable}{2}
\header{Part & Zulip channel}
\row{Overview, Motivation and history & overview}
\row{Techniques through examples & examples}
\row{Techniques and general results & techniques}
\row{Extensions and recent advances & advances}
\row{Toolkit & toolkit}
\row{Numerical methods & numerics}
\row{About this wiki & contributing}
\end{wikitable}
Discuss the mathematics, not the people. Associate editors moderate the channel and the talk pages of their part, and can hide messages; editors can remove a member whose participation disturbs the work rather than helping it. The Zulip is \href{https://endipa.zulipchat.com}{endipa.zulipchat.com}, by invitation; its administrators invite new members and create channels.

\section{Formulas in the wiki}
For editors who write the wiki's own format.
\begin{itemize}
  \item Inline formulas need \texttt{display="inline"}, displayed ones \texttt{display="block"}; each formula is written on one line. The renderer rejects \texttt{\textbackslash{}newcommand}, \texttt{\textbackslash{}label}, \texttt{\textbackslash{}eqref} and \texttt{\textbackslash{}substack}. No formulas in headings.
  \item Citations link to the Literature page by the key of the reference catalogue: \nowikicode{[[Literature\#\allowbreak{}BedrossianCotiZelati2017|\allowbreak{}Bedrossian and Coti Zelati (2017)]]}. The section \emph{References} of each page is generated from the catalogue: add the key to its comment line rather than editing the list.
  \item To link to a result, use the \texttt{id} of its box as the anchor, as in \texttt{Couette flow\#thm-couette}.
\end{itemize}

\section{The edition}
A static edition of this wiki, with the pictures, the videos and the TeX source of every page, is built from the wiki's text by a script and never edited by hand: \href{https://endipa.github.io}{endipa.github.io}.
\end{document}
